\documentclass[main.tex]{subfiles}
\begin{document}
\ifSubfilesClassLoaded{\begin{center}\small\itshape Part of the GnssProcessor Mathematical Reference \textbullet\ \today\end{center}}{}
\section{Estimator Infrastructure: WLS Seeding, State Initialisation, and Pivot Management}\label{sec:estinfra}
% ============================================================
%  Estimator-support machinery: bootstrap and re-reference
% ============================================================

% --- Doc-local symbols (descriptive names to avoid package collisions) ---
\newcommand{\wlsdesign}{\bm{A}}                 % WLS code-only design matrix
\newcommand{\wlsweight}{\bm{W}}                 % WLS diagonal weight matrix
\newcommand{\wlsstate}{\bm{\xi}}                % WLS core state [r, c dt, ZWD]
\newcommand{\mapw}{m_w}                          % wet mapping function
\newcommand{\coderes}{\widetilde{z}}            % prefit residual after modelled terms
\newcommand{\medop}{\operatorname{med}}         % median operator
\newcommand{\isbsys}{\mathrm{ISB}}              % inter-system bias seed
\newcommand{\pclk}{\delta_{\Phi}}              % phase-clock state
\newcommand{\ionoest}{\widehat{I}}              % per-sat slant iono estimate
\newcommand{\biascoef}{\beta}                   % receiver-bias Jacobian coefficient
\newcommand{\stransform}{\bm{T}}                % S-transform / re-referencing matrix
                                                % (distinct glyph: \Smat is the
                                                % innovation covariance elsewhere)

The Kalman filter of \dcref{sec:kf}{the filtering document} is a
\emph{local} linear estimator: every update assumes the state lies inside a
small neighbourhood of the linearisation point, and every dynamic
state---ambiguity, per-satellite ionosphere, receiver code bias, satellite
bias, inter-system offset (\dcref{sec:state}{the state document})---must be
born with a value and a covariance before it can be updated.  This document
covers the \emph{support} machinery that supplies those preconditions: the
weighted-least-squares (WLS) epoch solve that drops the receiver inside the
linearisation radius (\Cref{sec:estinfra-wls}); the first-epoch
residual-differencing seeds that give each newly created state a physically
informed prior (\Cref{sec:estinfra-init}); and the between-satellite pivot
selection and disturbance logic that keep the single-difference (SD)
reference frame stable across an arc (\Cref{sec:estinfra-pivot}).  None of
this is filtering proper---it is the scaffolding that makes the filter's
local-linear assumptions valid at the first epoch and after every
discontinuity.

The lossless covariance re-referencing algebra (the $\stransform\,\covM\,\stransform^\Tr$
S\nobreakdash-transform that re-anchors a group onto a new pivot) is
deliberately \emph{out of scope} here: it is derived in
\dcref{sec:diff}{the differencing document} and implemented once in
\texttt{include/filter/CovarianceReReference.hpp}, which both the
single-receiver and network handover paths call.  This section names that
seam where the pivot logic invokes it and otherwise refers the reader to the
differencing derivation.

% ============================================================
\subsection{The WLS epoch estimator: receiver-clock seed and post-gap re-seed}
\label{sec:estinfra-wls}
% ============================================================

Before the filter can run, the receiver position must be close enough to
truth that the geometry Jacobian \dcref{eq:geom-linear}{linearised about it}
is valid to sub-decimetre.  Broadcast pseudoranges alone fix this: a
GPS-only, code-only iterated Gauss--Newton solve over the compact core state
\begin{equation}
\wlsstate = \begin{bmatrix} \rxpos^\Tr & \sLight\,\mathrm{d}t_\rx & \zwd \end{bmatrix}^\Tr
\in\mathbb{R}^{5},
\label{eq:estinfra-wls-state}
\end{equation}
with $\rxpos\in\mathbb{R}^3$ the ECEF receiver position, $\sLight\,\mathrm{d}t_\rx$
the receiver clock in metres, and $\zwd$ the zenith wet delay
(\texttt{src/core/math/WlsEstimator.cpp}).  At least
\texttt{kCoreSize}~$=5$ code observations are required before the solve is
attempted; the iteration cap is \texttt{kMaxIter}~$=5$.

For each code measurement $\sat$ the prefit residual removes every modelled
term and the current state's contribution,
\begin{equation}
\coderes_\sat = \code_\sat
  - \bigl[\,\code_\sat^{\mathrm{mod}} + (\geomrange_\sat - \geomrange_\sat^{0})
  + \sLight\,\mathrm{d}t_\rx + \mapw(\sat)\,\zwd \,\bigr],
\label{eq:estinfra-wls-prefit}
\end{equation}
where $\code_\sat^{\mathrm{mod}}$ is the modelled pseudorange,
$\geomrange_\sat-\geomrange_\sat^{0}$ the change in geometric range relative
to the range at the iteration's linearisation point, and $\mapw(\sat)$ the
wet mapping function (\texttt{WlsEstimator::predictObs}, evaluated once per
row from the design-row loop of \texttt{WlsEstimator::estimate}).  The
geometric range is \emph{re-evaluated every iteration} about the running
position estimate inside that loop, so this is a genuine
Gauss--Newton relinearisation rather than a single-shot least-squares fit.

The design row for satellite $\sat$ is the negated line-of-sight stacked with
the clock and wet-mapping partials, assembled in the same loop:
\begin{equation}
\wlsdesign_\sat = \begin{bmatrix} -\los_\sat^\Tr & 1 & \mapw(\sat) \end{bmatrix}
= \begin{bmatrix} -e_x & -e_y & -e_z & 1 & \mapw(\sat) \end{bmatrix},
\label{eq:estinfra-wls-row}
\end{equation}
with $\los_\sat$ the receiver-to-satellite unit vector (the same line-of-sight
that drives \dcref{eq:los}{the observation geometry}).  Stacking all rows into
$\wlsdesign$, the prefit residuals into $\coderes$, and the inverse measurement
variances onto the diagonal weight $\wlsweight$ (clamped from below by
\texttt{kMinVariance}~$=10^{-6}$ in that same loop), the
normal-equations step is the boxed Gauss--Newton update
\begin{empheq}[box=\fbox]{equation}
\Delta\wlsstate = \bigl(\wlsdesign^\Tr \wlsweight\, \wlsdesign\bigr)^{-1}
                  \wlsdesign^\Tr \wlsweight\, \coderes,
\label{eq:estinfra-wls-step}
\end{empheq}
applied once per Gauss--Newton iteration of \texttt{WlsEstimator::estimate}.
Iteration terminates on a position
increment $\lVert\Delta\rxpos\rVert < \texttt{kTol}=0.1\,\mathrm{m}$
or at \texttt{kMaxIter}.  The normal matrix is
guarded against rank deficiency: if
\begin{equation}
\det\!\bigl(\wlsdesign^\Tr \wlsweight\, \wlsdesign\bigr) \le 10^{-6}
\label{eq:estinfra-wls-detguard}
\end{equation}
the solve aborts and the epoch is left un-seeded
(the determinant guard immediately preceding the step), which protects
against the degenerate
single-constellation, near-coplanar geometry that would otherwise produce a
wild position step.

On convergence the estimator returns both the state
$\wlsstate=[\rxpos^\Tr,\ \sLight\,\mathrm{d}t_\rx,\ \zwd]^\Tr$ and its formal
covariance, the inverse normal matrix
\begin{equation}
\covM_{\wlsstate} = \bigl(\wlsdesign^\Tr \wlsweight\, \wlsdesign\bigr)^{-1}
\label{eq:estinfra-wls-cov}
\end{equation}
(the result assembly at the tail of \texttt{WlsEstimator::estimate}).  Of
this pair the filter consumes only the
\emph{clock}: \texttt{Application::handleWlsSeeding} writes
$\sLight\,\mathrm{d}t_\rx$ into the receiver-clock state, and only when one of
three triggers fires---the first epoch, a data gap, or a clock \emph{jump} of at
least \texttt{kClockJumpThreshSec}~$=0.5\,\mathrm{ms}$ between the WLS and filter
clock estimates (in practice the jump, not the gap, is the dominant trigger).
The position, the zenith wet delay, and the covariance
\eqref{eq:estinfra-wls-cov} are computed but not consumed as filter seeds: the
position prior is the RINEX-header \texttt{approx\_position} plus the antenna
offset, and the WLS position re-seeds the filter only on a kinematic (non-IMU)
long-gap cold start (\texttt{Application::handleGap}); the ZWD estimate reaches
only a trace-CSV column; and the formal covariance has no filter consumer at
all.  The whole seeding path is skipped under BSSD, where the filter carries no
receiver-clock state for the WLS clock to seed
(\texttt{handleWlsSeeding} early-returns on \texttt{!hasClock()}).  The clock
seed's recurring role is \emph{post-gap re-seeding}: the receiver clock is the
most volatile state across a data gap (an unsteered oscillator can drift
arbitrarily while no measurements constrain it), so after a gap the WLS clock
estimate re-anchors $\sLight\,\mathrm{d}t_\rx$ before the phase filter
resumes, rather than letting the filter chase a stale clock prior through a
sequence of large innovations.

% ============================================================
\subsection{First-epoch state seeding by residual differencing}
\label{sec:estinfra-init}
% ============================================================

The WLS solve of \Cref{sec:estinfra-wls} seeds the receiver clock (and, on a
kinematic cold start, position); it does not touch the \emph{differential}
states.  Every such state---inter-system bias, phase clock, uncombined
receiver code bias, per-satellite ionosphere, carrier ambiguity---is created
empty and must be seeded from the measurement residuals.  The unifying
idea (\texttt{src/core/measurement/StateInitializer.cpp}) is that each such
state is precisely the quantity left in a particular \emph{difference} of
residuals once all already-modelled terms are removed; estimating it is a
matter of forming that difference and taking a robust (median) central value.
The shared prefit residual is
\begin{equation}
\coderes_{\rx,\sigl}^{\sat}
= \phaseo/\code_{\rx,\sigl}^{\sat}
  - \bigl(\mathrm{modelled}
  + \sLight\,\mathrm{d}t_\rx + \mapw\,\zwd + \isbsys + g_{\mathrm{trop}}
  + \alpha_\sigl \ionoest^\sat + \biascoef_\sigl b_\rx + b^\sat\bigr),
\label{eq:estinfra-rawres}
\end{equation}
the observation minus every term the model can already account for, evaluated
at the current state (\texttt{StateInitializer::computeRawResidual}, which
also forms the tropospheric gradient projection
$g_{\mathrm{trop}} = m_g(g_N\cos\mathrm{az} + g_E\sin\mathrm{az})$).  The
five \emph{differential} seeds (a)--(e) below are
differences of \eqref{eq:estinfra-rawres}; a sixth, network-only pre-seed~(f)
anchors the receiver clock itself.

\paragraph{(a) Inter-system bias.}  Each non-GPS constellation $g$ carries a
receiver clock offset relative to the GPS time scale
(\dcref{sec:state}{the inter-system offset state}).  Because the GPS clock is
the datum that the WLS step has already absorbed into
$\sLight\,\mathrm{d}t_\rx$, the ISB of $g$ is the constellation-median code
residual minus the GPS-median residual:
\begin{empheq}[box=\fbox]{equation}
\isbsys_g = \medop_{\sat\in g}\ \coderes^{\sat}
          - \medop_{\sat\in\mathrm{GPS}}\ \coderes^{\sat}
\label{eq:estinfra-isb}
\end{empheq}
(\texttt{StateInitializer::initializeInterSystemBiases}).  Medians, not means, are used throughout
so a single gross code outlier cannot capsize the seed.  A hard guard requires
at least one GPS code observation at the seeding epoch
(the GPS-median guard in that same routine): if GPS is absent it returns
\texttt{false} and \emph{retries the next epoch} rather than latching the ISB
states as initialised against an arbitrary datum.  This avoids the failure
mode where a constellation that happens to acquire first becomes the implicit
zero, and a GPS satellite rising later is then forced to wear the whole offset.

\paragraph{(b) Phase clock.}  In the uncombined model the phase observations
see a receiver clock that differs from the code clock by the receiver
phase--code hardware delay (\dcref{sec:state}{the phase-clock state}).  Its
seed is the median over satellites carrying \emph{both} a code and a phase
observation of the code-minus-phase residual:
\begin{equation}
\pclk = \medop_{\sat}\ \bigl(\coderes^{\sat}_{\code} - \coderes^{\sat}_{\phaseo}\bigr)
\label{eq:estinfra-pclk}
\end{equation}
(\texttt{StateInitializer::initializePhaseClockOffset}).  The seed's covariance
is then set from the code noise, tightened relative to the (looser) default
initial standard deviation but \emph{independent} of the $n$ differences that
formed the median (the covariance-seeding tail of that routine):
\begin{equation}
\sigma_{\pclk} = \max\!\Bigl(\tfrac{\sigma_{\code}}{\sqrt{n}},\ \sigma_{\code}\Bigr),
\qquad
\covM(\pclk,\pclk) = \sigma_{\pclk}^{2},
\label{eq:estinfra-pclk-cov}
\end{equation}
with the corresponding row and column of $\covM$ first zeroed so the seeded
state starts uncorrelated.  The $\max$ against $\sigma_{\code}$ makes the
$\sigma_{\code}/\sqrt{n}$ branch unreachable---$\sigma_{\code}/\sqrt{n}\le\sigma_{\code}$
for every $n\ge1$---so the seeded variance is $\sigma_{\code}^{2}$ regardless of
$n$: a conservative floor at the single-measurement code noise, not a
$1/\sqrt{n}$ promise.

\paragraph{(c) Uncombined receiver code bias / IFB.}  In the uncombined
PPP-RTK model each non-reference code band carries a receiver differential
code bias (DCB) or inter-frequency bias (IFB) state relative to the
constellation's BIA clock-reference band
(\dcref{sec:state}{the code-bias state}).  Forming the residual difference
between band $j$ and the reference band, removing the differential ionospheric
contribution (the two bands scale the same slant ionosphere by different
$\alpha$), and dividing by the difference of receiver-bias Jacobian
coefficients gives
\begin{empheq}[box=\fbox]{equation}
\widehat{b}_{\rx,j}
= \frac{\Delta\coderes_j - \Delta(\alpha_j\,\ionoest)}{\Delta\biascoef_j}
= \frac{(\coderes_j - \coderes_{\mathrm{ref}})
        - (\alpha_j - \alpha_{\mathrm{ref}})\,\ionoest^\sat}
       {\biascoef_j - \biascoef_{\mathrm{ref}}}
\label{eq:estinfra-codebias}
\end{empheq}
(\texttt{StateInitializer::initializeCodeBiases}), median-aggregated over satellites and
subject to the same $\lvert\widehat{b}\rvert > 200\,\mathrm{m}$ reject as the
ionosphere seed below.  The Jacobian coefficient $\biascoef_\sigl$ depends on
the bias parameterisation (\texttt{StateInitializer::resolveReceiverBias}):
for a pure
IFB state it is the unit coefficient $\biascoef = 1$, whereas for a
geometry-free / ionosphere-free DCB state against reference bands of
frequencies $f_1,f_2$ it is
\begin{equation}
\mu_1 = 1, \quad \mu_2 = \frac{f_1^2}{f_2^2}, \quad
\Delta\mu = \mu_1 - \mu_2, \qquad
\biascoef_j = \begin{cases} \mu_1/\Delta\mu & j = \text{band }1 \\[2pt]
                            \mu_2/\Delta\mu & j = \text{band }2, \end{cases}
\label{eq:estinfra-bias-coef}
\end{equation}
guarded by $\lvert\Delta\mu\rvert \ge 10^{-12}$ to avoid dividing by a
near-degenerate frequency separation.

A third, higher-priority branch precedes both cases for GLONASS: under
\texttt{enable\_glonass\_ifb} each FDMA band (G1, G2) carries one bias state
\emph{per FCN slot} $k$ with $\biascoef = 1$
(\texttt{StateKeys::makeGlonassPerSlotIfbKey}), because the per-slot carrier
$f = f_{\mathrm{base}} + k\,\Delta f$ gives a slot-dependent receiver group delay
that a single per-band state cannot represent; the central channel $k=0$ is the
held datum (no state) and the CDMA bands (G3/G4/G6) fall through to the legacy
per-band treatment above.  This branch runs before the clock-reference lookup
because GLONASS usually ships no clock-reference pair in the \texttt{.BIA}
product and would otherwise early-return before reaching any IFB logic; with the
flag off a non-reference GLONASS band carries no IFB state at all
(the per-slot and non-reference-band branches of
\texttt{resolveReceiverBias}).

\paragraph{(d) Per-satellite slant ionosphere.}  Each tracked satellite's
slant ionospheric delay (\dcref{sec:state}{the per-satellite iono state}) is
recovered from the geometry-free combination of two code bands.  With
ionospheric scale factors $\alpha_\sigl = f_{\mathrm{ref}}^2/f_\sigl^2$
referenced to GPS~L1 (\dcref{eq:iono-scaling}{the iono-scaling convention}),
the seed is
\begin{empheq}[box=\fbox]{equation}
\ionoest^\sat
= \frac{(\code_1 - \code_1^{\mathrm{mod}}) - (\code_2 - \code_2^{\mathrm{mod}})}
       {\alpha_1 - \alpha_2}
\label{eq:estinfra-iono}
\end{empheq}
(\texttt{StateInitializer::computeIonosphereSeed}), where $\code_k^{\mathrm{mod}}$ is the
modelled pseudorange on band $k$.  Two structural guards apply: the denominator
must satisfy $\lvert\alpha_1-\alpha_2\rvert \ge 10^{-6}$
(the denominator guard in that routine), and the seed is rejected unless it is
finite and $\lvert\ionoest^\sat\rvert < 200\,\mathrm{m}$
(its magnitude reject).  The $200\,\mathrm{m}$ reject is a physical
sanity bound---a slant ionospheric delay larger than this is never real and
signals a band mis-pairing or a gross code outlier, which must not be latched
into a per-satellite state.  The seed additionally applies two BIA-coverage
gates so the differenced code is bias-consistent: \emph{presence} (a
\texttt{.BIA} entry exists on the band, with the zero-valued BeiDou C2I/C6I
entries treated as missing), and---outside the network producer, which keeps
presence-only semantics---coverage \emph{honesty}
(\texttt{bias\_coverage\_requires\_apply}), which skips the seed when an entry is
loaded but not \emph{applied} to $(P-\mathrm{mod})$ so the raw hardware bias
still survives.  The honesty gate matters precisely because the failure it
blocks---a zero-valued BeiDou entry injecting $30$--$60\,\mathrm{m}$---sits well
under the $200\,\mathrm{m}$ reject above and would not be caught by it.

\paragraph{(e) Carrier ambiguity.}  The float ambiguity that subsequently feeds
integer resolution (\dcref{sec:ar}{the ambiguity-resolution document}) is what
remains of the \emph{phase} residual once every other modelled and estimated
term has been stripped (\texttt{StateManager::computeAmbiguitySeed}, the shared
routine every phase-seeding call site invokes),
\begin{equation}
\widehat{N}^\sat_{\rx,\sigl}
= \phaseo^\sat_{\rx,\sigl} - \mathrm{modelled}
  - \bigl(\sLight\,\mathrm{d}t_\rx + \pclk\bigr) - \mapw\,\zwd
  - \isbsys - g_{\mathrm{trop}} + \alpha_\sigl\,\ionoest^\sat,
\label{eq:estinfra-amb}
\end{equation}
i.e.\ the phase observation minus its model, receiver clock \emph{and} phase
clock $\pclk$ (seed~(b)), wet delay, inter-system bias, and tropospheric
gradient, plus the slant ionospheric phase advance $+\alpha_\sigl\,\ionoest^\sat$
(seed~(d))---the phase leads by the same $\alpha_\sigl$ the code lags, so the
iono enters with the sign opposite to \eqref{eq:estinfra-rawres}.  Both the
phase clock and the slant iono are removed here because they are already
estimated states in the uncombined model; the one term this subtraction
\emph{cannot} remove is a receiver clock still sitting at zero on a network cold
start, which is what pre-seed~(f) exists to anchor.  By construction this is the
phase residual after removing every term the filter already estimates, so the
seeded ambiguity
absorbs exactly the integer cycle count plus the residual receiver/satellite
phase-bias remainder---the quantity whose integer part AR recovers once the
satellite phase bias is supplied and the receiver phase bias cancels under
between-satellite differencing.

\paragraph{(f) Cold receiver-clock pre-seed (network path).}  Seeds (a)--(e)
assume the receiver clock has already been solved---true on the PPP path, where
the WLS clock seed of \Cref{sec:estinfra-wls} runs first.  The network producer
does not run that WLS seed, so on the network path nothing pre-solves the
receiver clock before the ambiguity of (e) is seeded:
\texttt{computeAmbiguitySeed} then subtracts a clock term of zero and the entire
free-running receiver-clock offset leaks uniformly into every cold ambiguity
seed.  The opt-in \texttt{seedColdReceiverClock}
(\texttt{network.cold\_clock\_seed}, default off) closes this by writing the
median GPS code residual---the very \texttt{gps\_median} that ISB seed (a)
computes and discards---into $x(\mathrm{clock})$.  It is self-limiting: it
overwrites only a \emph{cold} clock, treating any
$\lvert x(\mathrm{clock})\rvert > \texttt{kColdClockEpsilonM} = 10^{-6}\,\mathrm{m}$
as already solved and latching it untouched, and it returns \texttt{false}
(retry next epoch, do \emph{not} latch) when the epoch carries no GPS code
observation---mirroring the GPS-reference contract of seed (a)
(\texttt{StateInitializer::seedColdReceiverClock}).

% ============================================================
\subsection{BSSD pivot management}
\label{sec:estinfra-pivot}
% ============================================================

Between-satellite single differencing (BSSD, \dcref{sec:bssd}{the BSSD
document}) re-expresses each constellation's ambiguities relative to a
designated \emph{pivot} satellite $\piv{g}$, cancelling the receiver phase bias
and rendering the SD ambiguities integer.  The pivot is not arbitrary: a poor
choice degrades every SD row it anchors, and an unstable choice forces costly
re-referencing handovers.  \texttt{src/core/filter/PivotManager.cpp} selects and
maintains it; the baseline double-difference path uses a structurally analogous
\texttt{src/core/filter/DdPivotManager.cpp}.

\subsubsection{Lexicographic candidate ranking}
\label{sec:estinfra-pivot-rank}

Candidates are ordered \emph{lexicographically}: a hard signal-coverage tier
first, then a continuous score within the tier.  Coverage dominates because the
pivot must track every signal any non-pivot tracks---a non-pivot signal the
pivot does not carry is orphaned and dropped---so \emph{any} full-coverage
candidate outranks \emph{any} partial-coverage candidate regardless of score
(the lexicographic comparator in \texttt{PivotManager::buildEpochMaps}).
Concretely, candidate $a$ beats $b$ iff its
phase-signal count is higher, or the counts tie and its score is higher.

Within a coverage tier the tie-break score combines float-convergence and
geometry (\texttt{PivotManager::scorePivotCandidate}):
\begin{empheq}[box=\fbox]{equation}
\mathrm{score}(\sat) =
  -\min\!\bigl(1,\ \sqrt{\,v_{\max}(\sat)\,}\bigr)
  + \frac{\mathrm{el}(\sat)}{90}
  - \tfrac{1}{2}\,\mathbb{1}\!\left[\mathrm{el}(\sat) < 20^\circ\right],
\label{eq:estinfra-pivot-score}
\end{empheq}
where $v_{\max}(\sat)$ is the worst (largest) float-ambiguity variance among the
satellite's tracked phase signals and $\mathrm{el}(\sat)$ is its elevation in
degrees.  The first term penalises an un-converged pivot: on handover the new
pivot's variance is added to every same-group SD state (the diagonal of the
re-reference, $\covM'(i,i)=\covM(i,i)+\covM(\piv{g},\piv{g})-2\covM(i,\piv{g})$ ---
the diagonal of the \dcref{sec:diff}{S-transform}, implemented in
\texttt{include/filter/CovarianceReReference.hpp}),
so a pivot whose own state is still settling injects that uncertainty into the
entire constellation; the penalty saturates at $-1$ (one metre of equivalent
$1\sigma$ ambiguity uncertainty), with a flat $-1$ sentinel for a candidate that
has \emph{no} state yet (the negative-\texttt{max\_amb\_variance} sentinel
branch of that score).  On the first BSSD
epoch every candidate in a tier shares the sentinel, so the score collapses to
elevation alone and no special first-election branch is needed.  The
$\mathrm{el}/90$ term prefers high-elevation pivots, and the explicit
$-\tfrac12$ sub-$20^\circ$ penalty hard-discourages a low-elevation pivot whose
tropospheric and multipath residuals exceed the formal stochastic model
regardless of how converged its float state looks.  Both score terms only
ever \emph{rank}
candidates that have already cleared a hard elevation gate: a satellite below
\texttt{pivot\_min\_el} (\texttt{getBssdPivotMinElevation}) is barred from
election outright---in the freshest-peer scan, the online pick, and the
sticky-retention test, all three in
\texttt{PivotManager::selectPivotCandidates}---so the
$-\tfrac12$ penalty only breaks ties among survivors rather than being the
mechanism that keeps low satellites out.  The gate falls through deliberately
when \emph{no} satellite clears the mask (a genuine code-only-float
constellation): the pick reverts to the precomputed best/highest-elevation
satellite so the candidate set is never emptied (the fallback from
\texttt{pickExclusionAwareBest} to the precomputed \texttt{best\_sat} /
\texttt{highest\_sat}).  Sticky retention---the code's own
\emph{elevation hysteresis}---suppresses epoch-to-epoch flip-flop, but on that
same \emph{absolute} gate, not a challenger comparison: the incumbent is kept
whenever it still clears \texttt{pivot\_min\_el}, is not excluded, and stays
within the toe spread of \Cref{sec:estinfra-pivot-toe} (the sticky branch of
\texttt{selectPivotCandidates}), with no score evaluated on that branch, so
a strictly better challenger cannot displace a qualifying incumbent; the score
re-ranking runs only once the incumbent has \emph{failed} the gate.

\subsubsection{Ephemeris-freshness guardrail}
\label{sec:estinfra-pivot-toe}

A pivot anchored to a stale broadcast ephemeris steps \emph{every} SD row that
references it, because its extrapolated orbit and clock drift away from the
fresher peers it is differenced against (observed in production as a Galileo
pivot pinned to a one-hour-old I/NAV record that detonated the entire SD frame
while its peers were minutes fresh).  The guardrail is \emph{spread-gated}: a
candidate is barred from pivot election when its absolute ephemeris age
$\lvert t_{\mathrm{obs}} - t_{\mathrm{oe}}\rvert$ exceeds the constellation's
freshest qualified peer by more than
\begin{equation}
\texttt{kPivotToeSpreadMaxS} = 1200\ \mathrm{s} = 20\ \mathrm{min}
\label{eq:estinfra-pivot-toe}
\end{equation}
(the \texttt{freshestToeAgeAbs} scan in
\texttt{PivotManager::selectPivotCandidates}, applied both at the online
pick and at the sticky-retention test).  Gating on the \emph{spread} rather
than an
absolute age means that when every candidate is comparably fresh the ranking is
untouched---only an outlier-stale candidate is displaced.  The barred satellite
remains usable as a non-pivot; the freshest peer always passes (its difference
from itself is zero), so the candidate set is never emptied.

\subsubsection{Pivot disturbance detection and handover}
\label{sec:estinfra-pivot-dist}

By the BSSD construction a held pivot's own SD ambiguity is identically zero
with zero variance---it \emph{is} the reference, so $x(\piv{g}) = 0$ and
$\covM(\piv{g},\piv{g}) = 0$ exactly.  Any departure from that frozen state
signals that the filter has, through process noise or an update, perturbed the
reference and the SD frame is no longer self-consistent.
\texttt{isPivotDisturbed} flags this when, for any of the pivot's phase signals,
\begin{empheq}[box=\fbox]{equation}
\bigl\lvert x(\piv{g})\bigr\rvert > 10^{-6}
\quad\text{or}\quad
\covM(\piv{g},\piv{g}) > 10^{-12}
\label{eq:estinfra-pivot-disturbed}
\end{empheq}
(\texttt{PivotManager::isPivotDisturbed}).  The two thresholds are numerical-zero
tolerances, not tuned gates: $10^{-6}\,\mathrm{m}$ on the value and
$10^{-12}\,\mathrm{m}^2$ on the variance are tight enough to catch any genuine
leakage of state onto the pivot while tolerating floating-point round-off in the
re-reference arithmetic.  To prevent the predict step from \emph{manufacturing}
a spurious disturbance, the pivot's process noise is suppressed during
prediction so $\covM(\piv{g},\piv{g})$ does not grow off zero
(\dcref{sec:kf}{the freeze-state / $\Qmat$-suppression mechanics}); without that
suppression every epoch's $\Qmat$ would trip \eqref{eq:estinfra-pivot-disturbed}
and force a needless handover.

When a handover does fire---either because the incumbent was disturbed or because
ranking elected a different satellite---the outgoing pivot is classified.  A
\emph{clean} handover (the old pivot's state is invertible and frame-consistent)
can be re-referenced losslessly and remains smoother-inverse-referenceable, so the
RTS smoother can later undo it; a \emph{forced} or disturbed handover cannot be
inverted cleanly and marks a frame discontinuity.  In either case the actual
covariance re-anchoring---subtract the new pivot's row/column from each non-pivot
group member, then zero the new pivot's row/column---is performed in place by
\texttt{CovReReference::applyBssd} in
\texttt{include/filter/CovarianceReReference.hpp}, the single shared
implementation of the $\stransform\,\covM\,\stransform^\Tr$ transform
\dcref{sec:diff}{derived in the differencing document}; it is \emph{not}
re-derived here.  The baseline DD path differs structurally: rather than a
subtract-and-freeze on a single pivot it negates the new-pivot row/column and
subtracts from the others (a different S-transform), then re-keys the affected
ambiguity states via \texttt{StateManager::renameAmbiguity}
(\texttt{DdPivotManager::handleHandover}).  Both paths share only the post-transform
positive-definiteness guard (\texttt{enforcePsd}); the single-receiver user path
runs the transform unguarded, since at user scale the groups are small and
handovers rare enough that floating-point error in $\stransform\,\covM\,\stransform^\Tr$ does
not accumulate.

\ifSubfilesClassLoaded{\printbibliography}{}
\end{document}
